\documentclass[aspectratio=169]{beamer}
\usepackage[utf8]{inputenc}
\usepackage{graphicx}
\usepackage{booktabs}
\usepackage{xcolor}
\usepackage{tikz}

% Theme
\usetheme{Madrid}
\usecolortheme{default}
\setbeamertemplate{navigation symbols}{}

% Title Page Information
\title[Speech-to-Video Framework]{Multimodal Speech-to-Video Generation Framework for Object Learning in Specially Abled Education}
\author[Sandesh PY]{Sandesh PY}
\institute[VVCE, Mysuru]{
    Dept. of Computer Science \& Engineering \\
    Vidyavardhaka College of Engineering \\
    Mysuru, India
}
\date{}

\begin{document}

% Slide 1: Title
\begin{frame}
    \titlepage
\end{frame}

% Slide 2: Introduction
\begin{frame}{Introduction}
    \begin{itemize}
        \item \textbf{Goal}: Provide equal learning opportunities for specially abled learners, particularly hyperactive children.
        \item \textbf{Challenge}: Traditional educational tools (flashcards, static images) provide limited dynamic representation of objects and actions.
        \item \textbf{Solution}: A multimodal Speech-to-Video Generation Framework that converts spoken or typed object names into short educational videos.
        \item \textbf{Core Technologies}: Generative AI, Latent Diffusion Models, Natural Language Processing, and Speech Recognition.
    \end{itemize}
\end{frame}

% Slide 3: Proposed System Architecture
\begin{frame}{Proposed System Architecture}
    \begin{columns}
        \begin{column}{0.5\textwidth}
            \begin{itemize}
                \item \textbf{User Input}: Accepts speech or textual input.
                \item \textbf{Speech Recognition \& Translation}: Converts audio to text.
                \item \textbf{Prompt Generation}: Converts the object name into a structured template.
                \item \textbf{GPU-Based Video Generation}: Uses pretrained generative models.
            \end{itemize}
        \end{column}
        \begin{column}{0.5\textwidth}
            % Placeholder for Image 1 from paper
            \begin{figure}
                \centering
                \includegraphics[width=0.9\textwidth]{1.png}
                \caption{Proposed Speech-to-Video generation system architecture}
            \end{figure}
        \end{column}
    \end{columns}
\end{frame}

% Slide 4: Methodology \& Core Technologies
\begin{frame}{Methodology \& Core Technologies}
    \begin{itemize}
        \item \textbf{Speech Processing}: Automatic Speech Recognition (ASR) algorithms convert user audio signals into textual data.
        \item \textbf{Language Translation}: Neural Machine Translation models process regional script inputs (e.g., Kannada) and cross-translate them into English.
        \item \textbf{Prompt Engineering}: Natural Language Processing (NLP) uses predefined, structured templates to maintain consistency in the generated visual characteristics.
        \item \textbf{Video Generation Core}: Pretrained \textbf{Latent Diffusion Models} (and Generative Adversarial Networks) synthesize the final video frames mapped to the text prompt.
        \item \textbf{Hardware Acceleration}: The heavy inference pipeline of the generative deep learning models is executed independently on dedicated \textbf{GPU-enabled servers}.
    \end{itemize}
\end{frame}

% Slide 5: Multi-Scene Generation Methodology & Pipeline
\begin{frame}{Processing Pipeline \& Scene Composition}
    \vspace{-0.2cm}
    \begin{columns}
        \begin{column}{0.4\textwidth}
            \begin{block}{Why Multiple Scenes?}
                Instead of a single continuous shot, the system merges multiple small scenes.
            \end{block}
            \begin{itemize}
                \item Viewpoint changes.
                \item Object movement.
                \item Context/background changes.
            \end{itemize}
        \end{column}
        \begin{column}{0.6\textwidth}
            \begin{block}{Pipeline Execution}
                The scene composition involves a sequential pipeline:
            \end{block}
            \begin{itemize}
                \item \textbf{Step 1: Contextual Prompting} \\
                The core object is placed in diverse, educational environments to maintain attention.
                \item \textbf{Step 2: Keyframe Synthesis} \\
                The Diffusion model generates the base structural grids for each specific view.
                \item \textbf{Step 3: Temporal Smoothing} \\
                The individual scenes are merged sequentially. This provides an engaging, high-stimulus visual experience tailored for hyperactive learners.
            \end{itemize}
        \end{column}
    \end{columns}
\end{frame}

% Slide 6: Experimental Analysis & Results
\begin{frame}{Experimental Analysis \& Results}
    \begin{columns}
        \begin{column}{0.4\textwidth}
            \begin{itemize}
                \item Tested with objects: \textit{apple, dog, tree, flower}.
                \item Successfully generated 8-scene educational videos.
            \end{itemize}

            \vspace{0.3cm}
            \resizebox{\textwidth}{!}{%
            \begin{tabular}{@{}llc@{}}
            \toprule
            \textbf{Object} & \textbf{Scenes} & \textbf{Status} \\ \midrule
            Aeroplane       & 8 scenes        & Successful       \\
            Dog             & 8 scenes        & Successful       \\
            Lion            & 8 scenes        & Successful       \\ \bottomrule
            \end{tabular}%
            }
        \end{column}
        \begin{column}{0.6\textwidth}
            % Placeholder for Image 3 from paper
            \begin{figure}
                \centering
                \includegraphics[width=0.9\textwidth]{2.png}
                \caption{Multi-scene generation process showing eight frames}
            \end{figure}
        \end{column}
    \end{columns}
\end{frame}

% Slide 7: Conclusion
\begin{frame}{Conclusion}
    \begin{itemize}
        \item We presented a novel Speech-to-Video generation framework to assist specially abled learners via visual demonstrations.
        \item The prototype confirms that using pretrained generative models without requiring custom datasets is a viable approach.
        \item Dynamic visual representations effectively link speech or text-based concepts to clear visuals, bridging the learning gap.
        \item This highlights the broader potential of Generative AI in creating accessible educational materials.
    \end{itemize}
\end{frame}

% Slide 8: Questions
\begin{frame}[standout]
    \centering
    \Huge \textbf{Thank You!}\\
    \vspace{1cm}
    \Large Any Questions?
\end{frame}

\end{document}
